\documentclass[11pt,a4paper]{article}
\usepackage{webclases-lesson}
\lessonSetup{T0}{v0.8}{T0 - Herramientas para empezar Fisica}
\begin{document}
\shorthandoff{<>}
\color{Ink}
\coverHeader{T0 --- Herramientas para empezar F\'isica}{Unidades, vectores, medida y c\'alculo para dar sentido a un resultado.}
Antes de calcular conviene saber qu\'e buscamos, en qu\'e unidades debe salir y qu\'e
direcci\'on o tama\~no ser\'ia razonable.
T0 es la caja de herramientas de toda la F\'isica de Bachillerato: re\'une las
herramientas que reutilizar\'as en T1--T6 y en 2.\textordmasculine.
Las cajas \textbf{\textcolor{ToolLine}{Herramienta}} y
\textbf{\textcolor{DefLine}{Definici\'on}} marcan lo que debes tener siempre a mano.
\begin{lessonindex}
\indexentry{sec:b1}{1. Medir y comprobar unidades}
\indexentry{sec:b2}{2. Pensar con vectores}
\indexentry{sec:b3}{3. Cu\'anto confiar en una medida}
\indexentry{sec:cifras}{4. Cifras significativas}
\indexentry{sec:b4}{5. Leer cambios y acumulaciones}
\end{lessonindex}

\block{1}{Medir y comprobar unidades}{sec:b1}
\begin{defbox}{magnitud f\'isica}
Propiedad de un cuerpo o de un fen\'omeno que se puede medir y expresar como un
n\'umero acompa\~nado de una unidad.
\end{defbox}
Medir es comparar una magnitud con una unidad acordada. La longitud de una mesa no
cambia al expresarla como $1{,}20\units{m}$ o $120\units{cm}$: cambia el n\'umero porque
cambia la unidad. Una magnitud no es su unidad: la masa se mide en kilogramos.
\heading{Un lenguaje com\'un}
\begin{defbox}{Sistema Internacional de unidades (SI)}
Sistema de unidades acordado internacionalmente que fija siete unidades base,
a partir de las cuales se obtienen todas las dem\'as.
\end{defbox}
\begin{tabularx}{\linewidth}{@{}>{\raggedright\arraybackslash}p{4.6cm}Y@{}}
\toprule
Magnitudes base del SI & Unidades y s\'imbolos\\
\midrule
Longitud, masa, tiempo & metro (m), kilogramo (kg), segundo (s)\\
Corriente, temperatura & amperio (A), kelvin (K)\\
Cantidad de sustancia & mol (mol)\\
Intensidad luminosa & candela (cd)\\
\bottomrule
\end{tabularx}
Las magnitudes derivadas combinan las anteriores:
$[v]=LT^{-1}$, $[a]=LT^{-2}$ y $[F]=MLT^{-2}$.
Los corchetes indican \emph{dimensiones}, no valores num\'ericos.
Algunas unidades derivadas tienen nombre propio:

\begin{tabularx}{\linewidth}{@{}>{\raggedright\arraybackslash}p{3.3cm}>{\raggedright\arraybackslash}p{3.3cm}>{\raggedright\arraybackslash}p{3.2cm}Y@{}}
\toprule
Magnitud & Unidad (s\'imbolo) & Sale de & En unidades base SI\\
\midrule
Fuerza & newton (N) & $F=ma$ & $\mathrm{kg\,m/s^2}$\\
Trabajo y energ\'ia & julio (J) & $W=Fd$ & $\mathrm{kg\,m^2/s^2}$\\
Potencia & vatio (W) & $P=W/t$ & $\mathrm{kg\,m^2/s^3}$\\
Presi\'on & pascal (Pa) & $p=F/S$ & $\mathrm{kg/(m\,s^2)}$\\
\bottomrule
\end{tabularx}

\textbf{La unidad de un resultado.} Sustituye cada magnitud de la f\'ormula por su
unidad SI y simplifica. En $W=Fd$:
$\N\cdot\mathrm{m}=\mathrm{kg\,m/s^2}\cdot\mathrm{m}=\mathrm{kg\,m^2/s^2}=\mathrm{J}$.
Los n\'umeros puros, como el $\tfrac12$, no tienen unidad.

\textbf{Argumentos sin dimensi\'on.} Lo que va dentro de $e^{x}$, $\ln x$, $\sin x$ o
$\cos x$ debe ser un n\'umero sin dimensi\'on. En $\cos(\omega t)$, $[\omega]=T^{-1}$,
as\'i que $\omega t$ no tiene dimensi\'on y se expresa en radianes.
Un \'angulo en radianes es una raz\'on entre dos longitudes (arco entre radio):
el radi\'an es una unidad \guillemotleft entre comillas\guillemotright. En $e^{-t/\tau}$, $\tau$ es un tiempo.

\begin{toolbox}{an\'alisis dimensional}
\toolpart{Qu\'e hace} Comprueba si una f\'ormula es posible comparando las dimensiones
de sus t\'erminos.
\toolpart{F\'ormula} Si $A=B+C$, entonces $[A]=[B]=[C]$; los argumentos de
$e^{x}$, $\ln x$, $\sin x$ y $\cos x$ son adimensionales.
\toolpart{En F\'isica} T1 (ecuaciones del MRUA), T3 (unidades de fuerza),
T5 (trabajo, energ\'ia y potencia) y cualquier resultado final.
\toolpart{Ejemplo} $[\tfrac12at^2]=(LT^{-2})\,T^2=L$: puede ser una posici\'on.
\end{toolbox}

\begin{examplebox}{descartar una ecuaci\'on imposible}
\origin{Elaboraci\'on propia.}
\step{Planteamiento} Se propone $v=v_0+2xt$, con $x$ posici\'on y $t$ tiempo.
Decide si puede describir una velocidad; basta una justificaci\'on dimensional.
\solutionblock
\step{Modelo} Solo se pueden sumar magnitudes con las mismas dimensiones.
\step{C\'alculo} $[v_0]=L/T$, pero $[2xt]=LT$: la suma no es posible.
En cambio, $[at]=(L/T^2)T=L/T$, de modo que $v=v_0+at$ pasa esta prueba.
\step{Interpretaci\'on} La comprobaci\'on elimina errores, pero no demuestra una ley:
$v=v_0+2at$ tambi\'en es dimensionalmente compatible y no es el MRUA habitual.
\end{examplebox}

\begin{exercisebox}{deducir la unidad}
(a) La energ\'ia de un muelle es $E=\tfrac12kx^2$, con $x$ en metros. ¿En qu\'e unidad
se mide $k$? (b) ¿En qu\'e unidad sale la potencia $P=Fv$?
\solutionblock
(a) $k=2E/x^2$: $\mathrm{J/m^2}=\N\cdot\mathrm{m}/\mathrm{m^2}=\mathrm{N/m}=\mathrm{kg/s^2}$.
(b) $\N\cdot\mathrm{m/s}=\mathrm{J/s}=\mathrm{W}$, vatio.
\end{exercisebox}

\heading{Convertir es multiplicar por uno}
Los prefijos cambian la escala: kilo $10^3$, centi $10^{-2}$, mili $10^{-3}$ y
micro $10^{-6}$. Un factor como $1000\units{m}/1\units{km}$ vale uno.
Se orienta para cancelar la unidad que sobra:
\[
72{,}0\units{km/h}\,
\frac{1000\units{m}}{1\units{km}}\,
\frac{1\units{h}}{3600\units{s}}=20{,}0\mps.
\]
La conversi\'on exacta no limita las cifras significativas. En \'areas y vol\'umenes
se eleva \emph{todo} el factor: $1\units{cm^2}=(10^{-2}\units{m})^2=10^{-4}\units{m^2}$.

\begin{exercisebox}{una conversi\'on que cambia de escala}
Una pieza tiene volumen $250\units{cm^3}$ y masa $0{,}675\kg$.
Calcula su densidad en $\mathrm{kg/m^3}$ con tres cifras significativas.
\solutionblock
$V=250(10^{-2})^3=2{,}50\cdot10^{-4}\units{m^3}$.
Por tanto $\rho=m/V=2{,}70\cdot10^3\units{kg/m^3}$.
Usar $10^{-2}$ en lugar de $10^{-6}$ cambiar\'ia el volumen por un factor $10^4$.
La densidad debe tener dimensi\'on $ML^{-3}$.
\end{exercisebox}

\block{2}{Pensar con vectores}{sec:b2}
\begin{defbox}{escalar}
Magnitud que queda completamente determinada por un n\'umero y su unidad; por ejemplo,
la masa o la temperatura.
\end{defbox}
\begin{defbox}{vector}
Magnitud que necesita un m\'odulo, una direcci\'on y un sentido; por ejemplo,
el desplazamiento o la fuerza.
\end{defbox}
\begin{defbox}{m\'odulo de un vector}
Longitud de la flecha que representa a $\vec A$; se escribe $|\vec A|$ y nunca es negativo.
\end{defbox}
\textbf{Notaci\'on.} Toda letra con flecha, como $\vec A$ o $\Delta\vec r$, es un vector.
Una letra sin flecha es un escalar. El m\'odulo se escribe siempre $|\vec A|$.
Dos flechas paralelas pueden tener sentidos contrarios.
Como vectores libres, son iguales si coinciden m\'odulo, direcci\'on y sentido,
aunque los dibujemos en lugares distintos.

\heading{De la flecha a las componentes}
Un mismo vector puede escribirse de cuatro formas. Las tres primeras son par,
columna y combinaci\'on de vectores unitarios:
\[
\vec A=(4;\ 3)\units{m}
=\begin{pmatrix}4\\3\end{pmatrix}\units{m}
=(4\ii+3\jj)\units{m}.
\]
$\ii$ y $\jj$ tienen m\'odulo uno y se\~nalan los sentidos positivos de los ejes.
Las componentes son proyecciones con signo; no son las longitudes de dos fuerzas nuevas.
La cuarta forma, la polar, usa el m\'odulo y el \'angulo de la figura.
\begin{center}
\begin{tikzpicture}[x=1cm,y=1cm,font=\small]
\draw[axis] (-.2,0)--(5.2,0) node[right] {$x$ (m)};
\draw[axis] (0,-.2)--(0,3.8) node[above] {$y$ (m)};
\draw[guide] (4,0)--(4,3)--(0,3);
\draw[vector] (0,0)--(4,3) node[midway,sloped,above] {$r=|\vec A|=5\units{m}$};
\draw[vectorgreen] (0,-.35)--(4,-.35) node[midway,below] {$A_x=4\units{m}$};
\draw[vectoramber] (4.35,0)--(4.35,3) node[midway,right] {$A_y=3\units{m}$};
\draw[-{Latex[length=1.6mm]}] (.85,0) arc[start angle=0,end angle=36.869898,radius=.85];
\node at (1.3,.4) {$\theta$};
\node[Muted,font=\scriptsize] at (2,-1.15) {$\theta$ se mide desde $+x$, en sentido antihorario};
\end{tikzpicture}
\end{center}
En el tri\'angulo rect\'angulo, el coseno divide cateto adyacente entre hipotenusa;
el seno divide cateto opuesto entre hipotenusa, y la tangente divide
cateto opuesto entre adyacente.

\begin{toolbox}{descomposici\'on en componentes}
\toolpart{Qu\'e hace} Pasa de m\'odulo y \'angulo a componentes, y al rev\'es.
\toolpart{F\'ormula} \mbox{$A_x=|\vec A|\cos\theta$}, \mbox{$A_y=|\vec A|\sin\theta$};
al rev\'es, \mbox{$|\vec A|=\sqrt{A_x^2+A_y^2}$} y \mbox{$\tan\theta=A_y/A_x$}
(mirando los signos). El \'angulo $\theta$ se mide desde $+x$, en sentido antihorario.
\toolpart{En F\'isica} T2 (tiro oblicuo), T3 y T4 (peso en un plano inclinado),
T5 (trabajo de una fuerza inclinada).
\toolpart{Ejemplo} $|\vec v_0|=10\mps$ a $30^\circ$: $v_{0x}=8{,}66\mps$ y
$v_{0y}=5{,}00\mps$.
\end{toolbox}

\textbf{Forma polar.} Con $r=|\vec A|$ y $\theta$ medido desde $+x$ en sentido
antihorario, el vector se escribe
\[
\vec A=(r\cos\theta;\ r\sin\theta).
\]
En la figura, $\vec A=(5\cos36{,}9^\circ;\ 5\sin36{,}9^\circ)\units{m}=(4;\ 3)\units{m}$.
En 3D se a\~nade $A_z\hat k$ y $|\vec A|=\sqrt{A_x^2+A_y^2+A_z^2}$.

\begin{examplebox}{una fuerza inclinada y su cuadrante}
\origin{Elaboraci\'on propia sobre la descomposici\'on del original T0.}
\step{Planteamiento} Una fuerza $\vec F$ de m\'odulo $|\vec F|=20{,}0\N$ forma
$35^\circ$ con $+x$. Calcula sus componentes con tres cifras significativas.
\solutionblock
\step{Modelo} Ambas componentes ser\'an positivas. La horizontal es el cateto
adyacente al \'angulo dado.
\step{C\'alculo}
\[
F_x=20{,}0\cos35^\circ=16{,}4\N,\qquad
F_y=20{,}0\sin35^\circ=11{,}5\N.
\]
\step{Interpretaci\'on} Cada componente es menor que $|\vec F|$.
Con las cifras sin redondear, $\sqrt{F_x^2+F_y^2}=20{,}0\N$.
Sobre una mesa lisa, $F_x$ interviene en el movimiento horizontal;
el efecto vertical se estudiar\'a con Newton en T3.
\variation{A $145^\circ$, $F_x=-16{,}4\N$ y $F_y=+11{,}5\N$.
La arcotangente sola dar\'ia $-35^\circ$; los signos sit\'uan el vector en el segundo
cuadrante. Si $F_x=0$, la direcci\'on vertical se decide sin dividir por cero.}
\end{examplebox}

\heading{Vector unitario: solo la direcci\'on}
Un vector unitario tiene m\'odulo 1. Indica una direcci\'on y un sentido sin hablar
de tama\~no. Se obtiene dividiendo cada componente entre el m\'odulo.

\begin{toolbox}{vector unitario}
\toolpart{Qu\'e hace} Conserva la direcci\'on y el sentido de $\vec A\ne\vec 0$ y
cambia su m\'odulo a 1.
\toolpart{F\'ormula}
$\displaystyle\hat u_A=\frac{\vec A}{|\vec A|}
=\left(\frac{A_x}{\sqrt{A_x^2+A_y^2}};\ \frac{A_y}{\sqrt{A_x^2+A_y^2}}\right)$.
\toolpart{En F\'isica} T2 (direcci\'on de la velocidad), T3 (direcci\'on de una
fuerza, $\vec F=|\vec F|\,\hat u$) y 2.\textordmasculine\ (campos, $\hat u_r$).
\toolpart{Ejemplo} $(3;\ 4)/5=(0{,}6;\ 0{,}8)$.
\end{toolbox}

\begin{examplebox}{el unitario de $\vec A=(6;\ -8)$}
\origin{Elaboraci\'on propia.}
\step{Planteamiento} Halla el vector unitario en la direcci\'on de $\vec A=(6;\ -8)$
y comprueba que su m\'odulo es 1.
\solutionblock
\step{C\'alculo} $|\vec A|=\sqrt{6^2+(-8)^2}=\sqrt{100}=10$, luego
$\hat u_A=\left(\tfrac{6}{10};\ \tfrac{-8}{10}\right)=(0{,}6;\ -0{,}8)$.
\step{Comprobaci\'on} $|\hat u_A|^2=0{,}6^2+(-0{,}8)^2=0{,}36+0{,}64=1$.
Los signos se conservan: $\hat u_A$ apunta al cuarto cuadrante, como $\vec A$.
\step{Error t\'ipico} Dividir entre la suma de las componentes, $6+(-8)=-2$, da
$(-3;\ 4)$: su m\'odulo es 5 y adem\'as tiene el sentido contrario.
Se divide siempre entre el m\'odulo.
\end{examplebox}

\heading{Sumar, restar y cambiar la escala}
Para sumar gr\'aficamente, coloca la cola del segundo vector en la punta del primero:
la resultante une el inicio con el final. El paralelogramo expresa la misma suma.
Por componentes se operan por separado los ejes.

\begin{toolbox}{suma de vectores}
\toolpart{Qu\'e hace} Combina varios vectores en uno solo, la resultante.
\toolpart{F\'ormula} $\vec A+\vec B=(A_x+B_x;\ A_y+B_y)$,
$\vec A-\vec B=\vec A+(-\vec B)$ y $\lambda\vec A=(\lambda A_x;\ \lambda A_y)$.
\toolpart{En F\'isica} T1 y T2 (desplazamientos, composici\'on de velocidades),
T3 (fuerza resultante $\sum\vec F$), T6 (momento lineal total).
\toolpart{Ejemplo} $(4;\ 3)+(-1;\ 2)=(3;\ 5)$.
\end{toolbox}

Multiplicar por $\lambda$ cambia el m\'odulo por $|\lambda|$; si $\lambda<0$ invierte el
sentido. Con $\vec A=(4;\ 3)$ y $\vec B=(-1;\ 2)$ en metros:
$\vec A-\vec B=(5;\ 1)$ y $-2\vec B=(2;\ -4)$.

\heading{Dos productos con significados distintos}
El \textbf{producto escalar} mide alineaci\'on. Da un n\'umero con signo.
Es cero para vectores perpendiculares y negativo si forman un \'angulo mayor que $90^\circ$.

\begin{toolbox}{producto escalar}
\toolpart{Qu\'e hace} Mide cu\'anto apunta un vector en la direcci\'on de otro;
da un escalar.
\toolpart{F\'ormula}
$\vec A\cdot\vec B=|\vec A||\vec B|\cos\theta=A_xB_x+A_yB_y+A_zB_z$.
\toolpart{En F\'isica} T5 (trabajo, $W=\vec F\cdot\Delta\vec r$, y potencia,
$P=\vec F\cdot\vec v$); en 2.\textordmasculine, flujo de un campo.
\toolpart{Ejemplo} $(4;\ 3)\cdot(-1;\ 2)=-4+6=2$.
\end{toolbox}

El \textbf{producto vectorial} da un vector perpendicular al plano de los dos
vectores. El sentido sigue la mano derecha: al girar los dedos desde $\vec A$ hacia
$\vec B$ por el \'angulo menor, el pulgar se\~nala $\vec A\times\vec B$.
Invertir el orden invierte el resultado.

\begin{toolbox}{producto vectorial}
\toolpart{Qu\'e hace} Da un vector perpendicular a $\vec A$ y $\vec B$ cuyo m\'odulo
es el \'area del paralelogramo que forman.
\toolpart{F\'ormula} $|\vec A\times\vec B|=|\vec A||\vec B|\sin\theta$;
las componentes salen del determinante de abajo.
\toolpart{En F\'isica} En 2.\textordmasculine\ Bachillerato: momento de una fuerza,
momento angular y fuerza magn\'etica $\vec F=q\,\vec v\times\vec B$.
\toolpart{Ejemplo} $\ii\times\jj=\hat k$, pero $\jj\times\ii=-\hat k$.
\end{toolbox}

Si los dos vectores est\'an en el plano $xy$, solo queda la componente $\hat k$:
\[
(A_x;\ A_y;\ 0)\times(B_x;\ B_y;\ 0)=(A_xB_y-A_yB_x)\hat k.
\]
Con $\vec A=(2;\ 0;\ 0)\units{m}$ y $\vec B=(0;\ 3;\ 0)\units{m}$,
$\vec A\cdot\vec B=0\units{m^2}$, pero
$\vec A\times\vec B=6\hat k\units{m^2}$: no son operaciones intercambiables.
En ejes $x$ derecha, $y$ arriba, $+\hat k$ sale del papel ($\odot$);
$-\hat k$ entra ($\otimes$).

\heading{El producto vectorial en 3D: Sarrus y determinante}
En 3D las componentes salen de un determinante $3\times3$ cuya primera fila son los
vectores unitarios:
\[
\vec A\times\vec B=
\begin{vmatrix}\ii&\jj&\hat k\\A_x&A_y&A_z\\B_x&B_y&B_z\end{vmatrix}.
\]
\textbf{M\'etodo 1: regla de Sarrus.} Copia las dos primeras columnas a la derecha.
Suma los tres productos que bajan hacia la derecha y resta los tres que bajan hacia
la izquierda:
\begin{center}
\begin{tikzpicture}[x=1.15cm,y=.8cm,font=\small]
\foreach \c in {0,1,2} {
  \draw[GreenLine!80!black,thick] ($(\c,2)+(.18,-.22)$)--($(\c+2,0)+(-.18,.22)$);
}
\foreach \c in {2,3,4} {
  \draw[RedLine,thick,dashed] ($(\c,2)+(-.18,-.22)$)--($(\c-2,0)+(.18,.22)$);
}
\draw[thick] (-.4,-.3)--(-.4,2.3) (2.4,-.3)--(2.4,2.3);
\tikzset{entry/.style={fill=white,inner sep=1.5pt}}
\node[entry] at (0,2) {$\ii$}; \node[entry] at (1,2) {$\jj$}; \node[entry] at (2,2) {$\hat k$};
\node[entry,text=Muted] at (3,2) {$\ii$}; \node[entry,text=Muted] at (4,2) {$\jj$};
\node[entry] at (0,1) {$A_x$}; \node[entry] at (1,1) {$A_y$}; \node[entry] at (2,1) {$A_z$};
\node[entry,text=Muted] at (3,1) {$A_x$}; \node[entry,text=Muted] at (4,1) {$A_y$};
\node[entry] at (0,0) {$B_x$}; \node[entry] at (1,0) {$B_y$}; \node[entry] at (2,0) {$B_z$};
\node[entry,text=Muted] at (3,0) {$B_x$}; \node[entry,text=Muted] at (4,0) {$B_y$};
\node[right,align=left,font=\scriptsize] at (4.7,1)
  {\textcolor{GreenLine!70!black}{continua: bajan a la derecha, suman}\\[2pt]
   \textcolor{RedLine}{discontinua: bajan a la izquierda, restan}\\[2pt]
   \textcolor{Muted}{gris: columnas copiadas}};
\end{tikzpicture}
\end{center}
\[
\vec A\times\vec B=A_yB_z\,\ii+A_zB_x\,\jj+A_xB_y\,\hat k
-\left(A_yB_x\,\hat k+A_zB_y\,\ii+A_xB_z\,\jj\right).
\]
\textbf{M\'etodo 2: desarrollo por la primera fila.} Los signos alternan
$+\,-\,+$; cada determinante $2\times2$ vale
$\left|\begin{smallmatrix}a&b\\c&d\end{smallmatrix}\right|=ad-bc$:
\[
\vec A\times\vec B=
\ii\begin{vmatrix}A_y&A_z\\B_y&B_z\end{vmatrix}
-\jj\begin{vmatrix}A_x&A_z\\B_x&B_z\end{vmatrix}
+\hat k\begin{vmatrix}A_x&A_y\\B_x&B_y\end{vmatrix}.
\]
El error cl\'asico es olvidar el signo menos delante de $\jj$.

\begin{examplebox}{el mismo producto por dos caminos}
\origin{Elaboraci\'on propia.}
\step{Planteamiento} Calcula $\vec A\times\vec B$ para $\vec A=(2;\ 1;\ -1)$ y
$\vec B=(1;\ 3;\ 2)$ por Sarrus y por la primera fila. Comprueba el resultado.
\solutionblock
\step{Sarrus} Hacia la derecha: $(1)(2)\,\ii=2\ii$, $(-1)(1)\,\jj=-\jj$ y
$(2)(3)\,\hat k=6\hat k$. Hacia la izquierda: $(1)(1)\,\hat k=\hat k$,
$(-1)(3)\,\ii=-3\ii$ y $(2)(2)\,\jj=4\jj$.
\[
\vec A\times\vec B=(2\ii-\jj+6\hat k)-(\hat k-3\ii+4\jj)=5\ii-5\jj+5\hat k.
\]
\step{Primera fila}
\[
\vec A\times\vec B=\ii\,(1\cdot2-(-1)\cdot3)-\jj\,(2\cdot2-(-1)\cdot1)
+\hat k\,(2\cdot3-1\cdot1)=5\ii-5\jj+5\hat k.
\]
Los dos m\'etodos dan $\vec A\times\vec B=(5;\ -5;\ 5)$.
\step{Comprobaci\'on} El resultado debe ser perpendicular a $\vec A$:
$\vec A\cdot(\vec A\times\vec B)=2\cdot5+1\cdot(-5)+(-1)\cdot5=0$.
Si se olvida el menos de $\jj$, sale $(5;\ 5;\ 5)$ y el producto escalar con $\vec A$
da $10\ne0$: el error se detecta.
\end{examplebox}

\block{3}{Cu\'anto confiar en una medida}{sec:b3}
\begin{defbox}{precisi\'on}
Grado de concordancia entre medidas repetidas de la misma magnitud: cuanto menor es su
dispersi\'on, m\'as precisa es la medida.
\end{defbox}
\begin{defbox}{exactitud}
Grado de cercan\'ia entre el resultado de una medida, o la media de varias, y el valor
de referencia.
\end{defbox}
\begin{defbox}{error}
Diferencia entre el valor medido y el valor de referencia,
$e=x_{\mathrm{med}}-x_{\mathrm{ref}}$.
\end{defbox}
Otros dos conceptos completan la descripci\'on de una medida:

\begin{tabularx}{\linewidth}{@{}p{2.8cm}Y@{}}
\toprule
Concepto & Qu\'e informa\\
\midrule
Resoluci\'on & Menor cambio distinguible por el instrumento.\\
Incertidumbre & Duda cuantificada que acompa\~na al resultado.\\
\bottomrule
\end{tabularx}
Una balanza mal ajustada puede dar resultados muy agrupados y alejados de la
referencia: es precisa, pero no exacta, porque tiene un sesgo. Repetir medidas ayuda a
valorar la dispersi\'on; no corrige por s\'i solo una mala calibraci\'on.
La resoluci\'on tampoco equivale autom\'aticamente a la incertidumbre total.

\heading{Error, intervalo y porcentaje}
El error $e$ de la definici\'on lleva signo;
su m\'odulo es $|e|$ y el error relativo es $|e|/|x_{\mathrm{ref}}|$ cuando la
referencia no es cero. Medir $9{,}90\units{m}$ frente a $10{,}00\units{m}$ da
$e=-0{,}10\units{m}$ y error relativo del $1{,}00\,\text{\%}$.
Sin referencia no podemos conocer ese error: comunicamos $x\pm\Delta x$ y
el criterio usado para estimar $\Delta x$.
La incertidumbre relativa $\Delta x/|x|$ compara calidad entre escalas.

\begin{keybox}
\textbf{¿Respecto a qu\'e valor?} El \textbf{valor de referencia esperado,
te\'orico o tabulado} viene de fuera del experimento; por ejemplo,
$g=9{,}81\units{m/s^2}$. Si el problema lo da, el error y la exactitud se juzgan
respecto a \'el. La \textbf{media de las medidas $\bar x$} es la mejor estimaci\'on
que sale de nuestros datos, pero solo coincidir\'ia con el valor real si no hubiera
error sistem\'atico, algo que los datos solos no permiten saber.

\textbf{Regla:} si aparece el valor esperado, \'usalo como referencia para el error;
si no, usa la media y dilo. Con la media como referencia medimos
\textbf{dispersi\'on (precisi\'on)}, no exactitud. El resultado siempre se centra
en la medida, $\bar x\pm\Delta x$; el valor esperado sirve para calcular el error y
comprobar si cae dentro de ese intervalo.
\end{keybox}

\heading{Propagar una incertidumbre peque\~na}
Como estimaci\'on lineal conservadora, se suman incertidumbres absolutas al sumar o
restar, y relativas al multiplicar o dividir cantidades positivas:
\[
\Delta(a\pm b)\approx\Delta a+\Delta b,\qquad
\frac{\Delta(ab)}{|ab|}\approx\frac{\Delta a}{|a|}+\frac{\Delta b}{|b|},\qquad
\frac{\Delta(a/b)}{|a/b|}\approx\frac{\Delta a}{|a|}+\frac{\Delta b}{|b|}.
\]
Para $q=a^n$, $\Delta q/|q|\approx |n|\Delta a/|a|$.
Son aproximaciones de primer orden, v\'alidas para incertidumbres relativas peque\~nas
y lejos de divisiones por cero; no sustituyen un an\'alisis estad\'istico de correlaciones.

\begin{exercisebox}{rapidez con incertidumbre}
$t=(2{,}36\pm0{,}04)\units{s}$ y $d=(1{,}20\pm0{,}01)\units{m}$.
Estima $v=d/t$ con el criterio anterior; conserva dos cifras en $\Delta v$.
\solutionblock
\[
v=0{,}508474\ldots\mps,\qquad
\frac{\Delta v}{v}\approx\frac{0{,}01}{1{,}20}+\frac{0{,}04}{2{,}36}=0{,}02528\ldots
\]
$\Delta v\approx0{,}01285\mps$, luego $v=(0{,}508\pm0{,}013)\mps$.
Como contraste, los extremos
$1{,}19/2{,}40=0{,}4958$ y $1{,}21/2{,}32=0{,}5216\mps$ dan un intervalo parecido.
La contribuci\'on relativa del tiempo es la mayor; mejorar solo la regla ayudar\'ia menos.
\end{exercisebox}

\begin{examplebox}{el intervalo sobre una regla}
\origin{Elaboraci\'on propia.}
Cinco medidas en cm son $12{,}40;\ 12{,}41;\ 12{,}41;\ 12{,}43;\ 12{,}45$, con
resoluci\'on $0{,}01\units{cm}$. Su media es
$\bar x=62{,}10/5=12{,}42\units{cm}$ y las desviaciones son
$-0{,}02;\ -0{,}01;\ -0{,}01;\ +0{,}01;\ +0{,}03\units{cm}$.
Por tanto, la m\'axima desviaci\'on es $\Delta x=0{,}03\units{cm}$.
\begin{center}
\begin{tikzpicture}[x=1.15cm,y=1cm,font=\scriptsize]
\draw[axis] (-.35,0)--(9.45,0);
\foreach \i in {0,...,9} {
  \draw (\i,0)--(\i,-.14);
}
\foreach \i/\value in {0/{12{,}38},1/{12{,}39},2/{12{,}40},3/{12{,}41},4/{12{,}42},5/{12{,}43},6/{12{,}44},7/{12{,}45},8/{12{,}46},9/{12{,}47}} {
  \node[below=2pt] at (\i,0) {\value};
}
\draw[GreenLine!80!black,very thick] (1,1.02)--(7,1.02);
\draw[GreenLine!80!black,thick] (1,1.02)--(1,.82) (7,1.02)--(7,.82);
\node[above=3pt] at (4,1.02) {$[12{,}39;\ 12{,}45]$};
\fill[BrandBlue] (2,.25) circle (2.1pt);
\fill[BrandBlue] (3,.25) circle (2.1pt);
\fill[BrandBlue] (3,.42) circle (2.1pt);
\fill[BrandBlue] (5,.25) circle (2.1pt);
\fill[BrandBlue] (7,.25) circle (2.1pt);
\draw[AmberLine!92!black,thick] (4,.13)--(4,.68);
\node[above] at (4,.7) {$\bar x=12{,}42$};
\draw[RedLine,thick,dashed] (2,-.1)--(2,.68);
\node[below=15pt] at (2,0) {$L_{\rm ref}=12{,}40$};
\end{tikzpicture}
\end{center}
$L=(12{,}42\pm0{,}03)\units{cm}$ construye un intervalo que contiene todas las
medidas. Su incertidumbre relativa es $\Delta x/|\bar x|=0{,}03/12{,}42\approx0{,}24\,\text{\%}$.
Tambi\'en puede escribirse
$[x_{\min};\ x_{\max}]=[12{,}40;\ 12{,}45]\units{cm}$. En esta segunda notaci\'on,
$c=(x_{\min}+x_{\max})/2=12{,}425$ y
$\Delta=(x_{\max}-x_{\min})/2=0{,}025\units{cm}$, la semidispersi\'on.
Los corchetes $[a;\ b]$ y $c\pm\Delta$ son dos formas de nombrar un mismo
intervalo; en F\'isica suele preferirse $\pm$. Centrar en la media y usar la
m\'axima desviaci\'on da un intervalo algo m\'as ancho que sigue conteniendo
todas las medidas.
\end{examplebox}

\block{4}{Cifras significativas}{sec:cifras}
Una cifra significativa es un d\'igito que aporta informaci\'on sobre la medida:
comunica hasta qu\'e decimal defendemos su precisi\'on. Los ceros iniciales solo
sit\'uan la coma; los dem\'as pueden comunicar precisi\'on.
\begin{tabularx}{\linewidth}{@{}p{5.2cm}Y@{}}
\toprule
Regla & Ejemplo\\
\midrule
D\'igitos distintos de cero cuentan; los ceros entre ellos tambi\'en. &
$305$ tiene 3 cifras significativas.\\
Ceros iniciales no cuentan. & $0{,}00450$ tiene 3.\\
Ceros finales tras la coma cuentan. & $12{,}30$ tiene 4.\\
Ceros finales de un entero son ambiguos: usa notaci\'on cient\'ifica. &
$1200$: $1{,}2\cdot10^3$ tiene 2; $1{,}200\cdot10^3$ tiene 4.\\
Conteos exactos y conversiones definidas no limitan las cifras. &
$1\units{m}=100\units{cm}$ es exacta.\\
\bottomrule
\end{tabularx}

\heading{Redondear y operar}
$0{,}004587$ a 2 cifras significativas es $0{,}0046$; $2{,}349$ a 2 es $2{,}3$.
En sumas y restas manda el dato con menos decimales:
$12{,}3+0{,}45=12{,}75\to12{,}8$ y
$105{,}3+2{,}456=107{,}756\to107{,}8$. En productos y cocientes manda el menor
n\'umero de cifras significativas:
$2{,}4\cdot3{,}12=7{,}488\to7{,}5$ y
$3{,}0\cdot4{,}56/1{,}2345=11{,}08\ldots\to11$.
Conserva cifras de guarda durante el c\'alculo y redondea solo al final.

\heading{Cuando hay una incertidumbre expl\'icita}
Redondea primero $\Delta x$ a una cifra (o a dos si empieza por 1) y despu\'es
el valor hasta el mismo decimal:
$9{,}8137\pm0{,}0462\to(9{,}81\pm0{,}05)$. El ejemplo de la regla se escribe,
con el mismo criterio, como $(12{,}42\pm0{,}03)\units{cm}$.

\begin{exercisebox}{contar y redondear}
Indica las cifras significativas de $0{,}0305$, $250{,}0$ y $2{,}50\cdot10^2$.
Redondea $4{,}376\pm0{,}084$ siguiendo la regla de incertidumbres.
\solutionblock
$0{,}0305$ tiene 3; $250{,}0$ tiene 4; $2{,}50\cdot10^2$ tiene 3.
La incertidumbre empieza por 8: $0{,}084\to0{,}08$, y el valor al mismo decimal es
$4{,}38$. Resultado: $(4{,}38\pm0{,}08)$.
\end{exercisebox}

\block{5}{Leer cambios y acumulaciones}{sec:b4}
Una gr\'afica temporal no dibuja el camino en el espacio: el eje horizontal es el
tiempo. La pendiente compara el cambio vertical con el horizontal y conserva unidades.
En $x(t)$, $\Delta x/\Delta t$ es una velocidad media; la derivada da la instant\'anea.

\heading{De la velocidad media a la instant\'anea}
Sea $x(t)=t^2/4$, con $x$ en metros y $t$ en segundos. Queremos la velocidad en
$t=2\units{s}$, que todav\'ia no sabemos calcular. S\'i sabemos calcular una media:
$v_{\mathrm{med}}=\Delta x/\Delta t$. Como $x(2)=1\units{m}$, tomamos intervalos que
empiezan en $2\units{s}$ y los hacemos cada vez m\'as cortos.

\begin{minipage}[c]{.36\linewidth}
\centering\small
\begin{tabular}{@{}lcc@{}}
\toprule
Intervalo & $\Delta t$ (s) & $v_{\mathrm{med}}$ (m/s)\\
\midrule
$[2;\ 4]$ & $2$ & $1{,}5$\\
$[2;\ 3]$ & $1$ & $1{,}25$\\
$[2;\ 2{,}5]$ & $0{,}5$ & $1{,}125$\\
$[2;\ 2{,}25]$ & $0{,}25$ & $1{,}0625$\\
$[2;\ 2{,}1]$ & $0{,}1$ & $1{,}025$\\
$[2;\ 2{,}01]$ & $0{,}01$ & $1{,}0025$\\
\bottomrule
\end{tabular}
\end{minipage}\hfill
\begin{minipage}[c]{.62\linewidth}
\centering
\begin{tikzpicture}[x=1.25cm,y=.95cm,font=\small]
\draw[axis] (0,0)--(4.75,0) node[right] {$t$ (s)};
\draw[axis] (0,0)--(0,5.2) node[above] {$x$ (m)};
\foreach \t in {1,2,3,4} \draw (\t,0)--(\t,-.08) node[below] {$\t$};
\foreach \p in {1,2,3,4} \draw (0,\p)--(-.06,\p) node[left] {$\p$};
\draw[guide] (2,0)--(2,1)--(0,1);
\draw[AmberLine!45,thick] (2,1)--(4.6,4.9);
\draw[AmberLine!85!black,thick] (2,1)--(4.6,4.25);
\draw[AmberLine!55!black,thick] (2,1)--(4.6,3.925);
\draw[GreenLine!80!black,dashed,thick] (1,0)--(4.6,3.6);
\draw[BrandBlue,thick,domain=0:4.2,samples=60] plot(\x,{\x*\x/4});
\foreach \t in {2.5,3,4} \fill[AmberLine!70!black] (\t,{\t*\t/4}) circle (1.6pt);
\fill[BrandBlue] (2,1) circle (2pt) node[above left] {$(2;\ 1)$};
\node[right,font=\scriptsize,AmberLine!70!black] at (4.6,4.9) {pendiente $1{,}5$};
\node[right,font=\scriptsize,AmberLine!85!black] at (4.6,4.25) [yshift=1.5pt] {$1{,}25$};
\node[right,font=\scriptsize,AmberLine!55!black] at (4.6,3.925) [yshift=-1pt] {$1{,}125$};
\node[right,font=\scriptsize,GreenLine!70!black] at (4.6,3.45) [yshift=-2pt] {tangente: $1$};
\end{tikzpicture}
\end{minipage}

\smallskip
Cada vez que el intervalo se hace m\'as peque\~no, la media se acerca m\'as a
$1\mps$. En la gr\'afica, cada secante (naranja, cada vez m\'as oscura) se acerca a la
recta tangente en $(2;\ 1)$, discontinua.

\textbf{El paso general.} Con un intervalo de duraci\'on $h$:
\[
v_{\mathrm{med}}=\frac{x(2+h)-x(2)}{h}=\frac{(2+h)^2/4-1}{h}=\frac{4h+h^2}{4h}=1+\frac h4 .
\]
Cuando $h$ se hace infinitamente peque\~no (un intervalo infinitesimal),
\[
v(2)=\lim_{h\to0}\left(1+\frac h4\right)=1\mps.
\]
¿Por qu\'e un l\'imite? Un solo instante tiene $\Delta t=0$, y no se puede dividir
entre cero. Pero s\'i podemos ver hacia d\'onde se dirigen las medias.

\begin{defbox}{derivada}
La derivada de $x(t)$ en $t$ es el l\'imite de $[x(t+h)-x(t)]/h$ cuando $h\to0$;
gr\'aficamente, es la pendiente de la recta tangente en ese punto.
\end{defbox}

\heading{Reglas m\'inimas de c\'alculo}
Una constante deriva a cero; $\frac{d}{dt}(ct^n)=cnt^{n-1}$.
La derivada de una suma es la suma de las derivadas.
\textbf{Comprobaci\'on:} $\frac{d}{dt}(t^2/4)=t/2$, y en $t=2\units{s}$ da
$1\mps$, lo mismo que la tabla.
Otro ejemplo: $x(t)=2+3t+t^2$ da $v=0+3+2t$ y $a=dv/dt=2\mpss$.
Los coeficientes llevan las unidades necesarias: en $t^2$ el coeficiente es
$1\units{m/s^2}$, no un n\'umero sin unidades.

\begin{toolbox}{derivada}
\toolpart{Qu\'e hace} Da el ritmo de cambio instant\'aneo de una magnitud.
\toolpart{F\'ormula} $\frac{d}{dt}(ct^n)=cnt^{n-1}$; $v=dx/dt$ y $a=dv/dt$.
\toolpart{En F\'isica} T1 y T2 (de $x$ a $v$ y de $v$ a $a$) y m\'aximos y m\'inimos,
como la altura m\'axima o el punto de retorno.
\toolpart{Ejemplo} $x=t^2/4\ \Rightarrow\ v=t/2$.
\end{toolbox}

\heading{Una barra de futbol\'in: avanzar, retroceder y chutar}
La barra de un futbol\'in se mueve en una sola dimensi\'on, a lo largo de su eje.
El jugador empuja hacia delante, retrocede un poco y despu\'es chuta. Un modelo es
\[
x(t)=2t^3-9t^2+12t\quad(x\text{ en cm},\ t\text{ en s},\ 0\le t\le3),
\]
con $x(0)=0$, $x(1)=5$, $x(2)=4$ y $x(3)=9\units{cm}$. Su velocidad es
\[
v(t)=6t^2-18t+12=6(t-1)(t-2)\units{cm/s}.
\]

\begin{minipage}[c]{.4\linewidth}
\centering\small
\begin{tabular}{@{}lcc@{}}
\toprule
Tramo (s) & Signo de $v$ & La barra\\
\midrule
$(0;\ 1)$ & $+$ & avanza\\
$t=1$ & $0$ & se para\\
$(1;\ 2)$ & $-$ & retrocede\\
$t=2$ & $0$ & se para\\
$(2;\ 3)$ & $+$ & avanza\\
\bottomrule
\end{tabular}
\end{minipage}\hfill
\begin{minipage}[c]{.58\linewidth}
\centering
\begin{tikzpicture}[x=2cm,y=.42cm,font=\small]
\draw[axis] (0,0)--(3.35,0) node[right] {$t$ (s)};
\draw[axis] (0,0)--(0,10) node[above] {$x$ (cm)};
\foreach \t in {1,2,3} \draw (\t,0)--(\t,-.2) node[below] {$\t$};
\foreach \p in {4,5,9} \draw (0,\p)--(-.04,\p) node[left] {$\p$};
\draw[guide] (1,0)--(1,5)--(0,5) (2,0)--(2,4)--(0,4) (3,0)--(3,9)--(0,9);
\draw[BrandBlue,thick,domain=0:3,samples=80] plot(\x,{2*\x^3-9*\x^2+12*\x});
\draw[GreenLine!80!black,thick] (.6,5)--(1.4,5);
\draw[GreenLine!80!black,thick] (1.6,4)--(2.4,4);
\fill[BrandBlue] (1,5) circle (2pt) node[above=3pt,font=\scriptsize] {m\'aximo, $v=0$};
\fill[BrandBlue] (2,4) circle (2pt) node[below=4pt,font=\scriptsize,fill=white,inner sep=1pt] {m\'inimo, $v=0$};
\node[font=\scriptsize,Muted] at (.5,1) {avanza $\rightarrow$};
\node[font=\scriptsize,Muted] at (1.5,1) {$\leftarrow$ retrocede};
\node[font=\scriptsize,Muted] at (2.5,1) {avanza $\rightarrow$};
\end{tikzpicture}
\end{minipage}

\smallskip
\textbf{La l\'ogica.} Para cambiar de sentido, la barra tiene que pararse un instante:
ah\'i $v=0$. \textbf{La definici\'on lo confirma.} Cerca de $t=1\units{s}$, las
secantes pasan de subir a bajar: entre $0{,}9$ y $1\units{s}$ la pendiente media es
$+0{,}32\units{cm/s}$, y entre $1$ y $1{,}1\units{s}$ es $-0{,}28\units{cm/s}$.
La tangente, que queda entre ambas, es horizontal: $v(1)=0$. Ah\'i la posici\'on es
m\'axima ($5\units{cm}$). En $t=2\units{s}$ ocurre lo contrario: m\'inimo de $4\units{cm}$.

\heading{M\'aximos y m\'inimos}
En un m\'aximo o un m\'inimo la derivada vale cero \textbf{y cambia de signo}.
Si pasa de $+$ a $-$, hay un m\'aximo; si pasa de $-$ a $+$, hay un m\'inimo.

\begin{warnbox}
$f'=0$ no siempre indica un extremo. $f(x)=x^3$ tiene $f'(x)=3x^2$, que vale cero en
$x=0$; pero $f'\ge0$ a ambos lados, no cambia de signo y $x^3$ sigue creciendo.
\end{warnbox}

\begin{exercisebox}{buscar y clasificar extremos}
Para $f(x)=x^3-3x^2$, halla d\'onde $f'=0$ y clasifica cada punto.
\solutionblock
$f'(x)=3x^2-6x=3x(x-2)$, que se anula en $x=0$ y $x=2$.
$f'$ pasa de $+$ a $-$ en $x=0$: m\'aximo, con $f(0)=0$.
Pasa de $-$ a $+$ en $x=2$: m\'inimo, con $f(2)=8-12=-4$.
\end{exercisebox}

\heading{Acumular: la integral}
Integrar suma peque\~nas contribuciones: velocidad por intervalo de tiempo da
desplazamiento. Una primitiva de $ct^n$ es $ct^{n+1}/(n+1)$ si $n\ne-1$.
La integral definida se calcula evaluando al final \emph{menos} al principio:
\[
\int_{t_1}^{t_2}v(t)\,dt=[V(t)]_{t_1}^{t_2}=V(t_2)-V(t_1).
\]
La integral indefinida incluye una constante $C$; una condici\'on inicial la fija.
Para $v=3+2t$, una posici\'on posible es $3t+t^2+C$.
Si $x(0)=2\units{m}$, entonces $C=2\units{m}$.

\begin{toolbox}{integral}
\toolpart{Qu\'e hace} Acumula peque\~nas contribuciones: es el \'area bajo la curva.
\toolpart{F\'ormula} $\int_{t_1}^{t_2}v\,dt=\Delta x$; una primitiva de $ct^n$ es
$ct^{n+1}/(n+1)$.
\toolpart{En F\'isica} T1 (de $v$ a $x$) y T5 (trabajo de una fuerza variable,
$W=\int F\,dx$).
\toolpart{Ejemplo} Con $v=t$ (SI) entre $0$ y $2\units{s}$: $\Delta x=[t^2/2]_0^2=2\units{m}$.
\end{toolbox}

\begin{examplebox}{integrar y recuperar la posici\'on}
\origin{Elaboraci\'on propia; continuaci\'on del mismo movimiento.}
\step{Planteamiento} Calcula el desplazamiento entre $0$ y $2\units{s}$ para
$v(t)=3+2t$ y comprueba $x(2)$.
\solutionblock
\step{Modelo} En $v(t)$, el \'area es un rect\'angulo m\'as un tri\'angulo.
\step{C\'alculo}
\[
\Delta x=\int_0^2(3+2t)\,dt=[3t+t^2]_0^2=10\units{m}.
\]
Geom\'etricamente: $3\cdot2+\tfrac12(2)(7-3)=10\units{m}$.
\step{Interpretaci\'on} La integral no entrega por s\'i sola la posici\'on:
$x(2)=x(0)+\Delta x=2+10=12\units{m}$.
En $a(t)$, el \'area entrega $\Delta v$, en $\mathrm{m/s}$.
\end{examplebox}

\begin{exercisebox}{\'area con signo}
Una velocidad vale $+2{,}00\mps$ durante $3{,}00\units{s}$ y luego
$-1{,}00\mps$ durante $2{,}00\units{s}$. Calcula desplazamiento y distancia con tres
cifras significativas. El cambio instant\'aneo se idealiza al separar los tramos.
\solutionblock
$\Delta x=+6{,}00-2{,}00=+4{,}00\units{m}$, mientras $d=6{,}00+2{,}00=8{,}00\units{m}$.
El \'area negativa corresponde a volver, no a una distancia negativa.
T1 conectar\'a esta lectura con posiciones, velocidades y aceleraciones del mismo m\'ovil.
\end{exercisebox}

\heading{Fuentes y procedencia}
Redacci\'on, ejemplos num\'ericos y diagramas propios a partir del alcance del
original T0 y su background editorial. El material de intake conserva unidades,
vectores, medida y c\'alculo; no se reproducen sus capturas.
Referencia de inspiraci\'on principal:
\source{https://fisquiweb.es/Apuntes/apun1B.htm}{FisQuiWeb / IES La Magdalena: apuntes de 1.\textordmasculine\ Bachillerato.}
Contraste externo: distinguir exactitud, precisi\'on y cifras comunicadas; interpretar
geom\'etricamente los dos productos vectoriales.
\source{https://openstax.org/books/university-physics-volume-1/pages/1-6-significant-figures}{OpenStax, University Physics 1, 1.6: Significant Figures.}
\source{https://openstax.org/books/calculus-volume-3/pages/2-3-the-dot-product}{OpenStax, Calculus 3, 2.3: The Dot Product.}
\source{https://openstax.org/books/calculus-volume-3/pages/2-4-the-cross-product}{OpenStax, Calculus 3, 2.4: The Cross Product.}
\end{document}
